% Standalone problem document, generated from the adjacent README.md.
% Compile with XeLaTeX. Mathematical content is maintained in Markdown.
\documentclass[11pt,a4paper]{article}
\usepackage[fontsize=10.5pt]{scrextend}
\usepackage[margin=22mm,headheight=15pt,headsep=9mm,footskip=11mm]{geometry}
\usepackage{fontspec}
\setmainfont{texgyrepagella}[Extension=.otf,UprightFont=*-regular,BoldFont=*-bold,ItalicFont=*-italic,BoldItalicFont=*-bolditalic]
\setsansfont{texgyreheros}[Extension=.otf,UprightFont=*-regular,BoldFont=*-bold,ItalicFont=*-italic,BoldItalicFont=*-bolditalic]
\setmonofont{texgyrecursor}[Extension=.otf,UprightFont=*-regular,BoldFont=*-bold,ItalicFont=*-italic,BoldItalicFont=*-bolditalic,Scale=MatchLowercase]
\usepackage{amsmath,amssymb}
\usepackage{unicode-math}
\setmathfont{texgyrepagella-math.otf}
\usepackage{xcolor,microtype,enumitem,needspace,fancyhdr}
\usepackage{bookmark}
\definecolor{ink}{HTML}{17344B}
\definecolor{accent}{HTML}{197D80}
\definecolor{muted}{HTML}{596773}
\hypersetup{colorlinks=true,linkcolor=accent,urlcolor=accent,pdftitle={MD-05 - The
sharp universal Spencer discrepancy
constant},pdfauthor={Open Problems in Numerical Linear Algebra}}
\urlstyle{same}
\setlength{\parindent}{0pt}
\setlength{\parskip}{5pt plus 1pt minus 1pt}
\linespread{1.04}
\setlength{\emergencystretch}{3em}
\widowpenalty=10000
\clubpenalty=10000
\displaywidowpenalty=10000
\allowdisplaybreaks[1]
\setlist{nosep,leftmargin=1.5em}
\providecommand{\tightlist}{\setlength{\itemsep}{0pt}\setlength{\parskip}{0pt}}
\providecommand{\pandocbounded}[1]{#1}
\pagestyle{fancy}
\fancyhf{}
\fancyhead[L]{\sffamily\footnotesize\color{muted}OPEN PROBLEMS IN NUMERICAL LINEAR ALGEBRA}
\fancyhead[R]{\sffamily\bfseries\footnotesize\color{accent}MD-05}
\fancyfoot[L]{\parbox[b]{0.9\textwidth}{\sffamily\fontsize{8}{10}\selectfont\color{muted}Editorial ratings; a bounded search cannot certify that a problem remains unsolved.\\Literature check: 2026-09-10}}
\fancyfoot[R]{\sffamily\footnotesize\color{muted}\thepage}
\renewcommand{\headrulewidth}{0.3pt}
\renewcommand{\headrule}{\hbox to\headwidth{\color{accent}\leaders\hrule height \headrulewidth\hfill}}
\makeatletter
\renewcommand{\section}{\Needspace{5\baselineskip}\@startsection{section}{1}{0pt}{12pt plus 2pt minus 2pt}{4pt}{\sffamily\bfseries\large\color{ink}}}
\renewcommand{\subsection}{\Needspace{4\baselineskip}\@startsection{subsection}{2}{0pt}{9pt plus 1pt minus 1pt}{3pt}{\sffamily\bfseries\normalsize\color{ink}}}
\makeatother
\setcounter{secnumdepth}{0}
\begin{document}
{\sffamily\small\color{muted}Matrix discrepancy and optimization\par}
\vspace{3pt}
{\raggedright\hyphenpenalty=10000\exhyphenpenalty=10000\sffamily\bfseries\fontsize{21}{24}\selectfont\color{ink}The
sharp universal Spencer discrepancy constant\par}
\vspace{7pt}
{\sffamily\small\textbf{Difficulty:} extreme\quad\textbf{Importance:} interesting
to the community\par}
{\sffamily\small\bfseries\color{accent}Status: Open\par}
\vspace{4pt}
{\color{accent}\hrule height 0.8pt}
\vspace{6pt}
\textbf{Provenance:} source-stated extremal-value problem.

\textbf{Rating rationale:} Determining the best constant across all sign
matrices and all orders is a longstanding extremal discrepancy barrier;
it sharpens a central matrix-balancing theorem.

For each positive integer \(n\), define \[
d_n=\frac{1}{\sqrt n}
\max_{A\in\{-1,1\}^{n\times n}}
\min_{x\in\{-1,1\}^n}\|Ax\|_\infty.
\] Determine the exact value of \[
C_*=\sup_{n\geq1}d_n.
\] Thus \(C_*\) is the smallest real constant for which every square
sign matrix \(A\) admits a column signing \(x\) with
\(\|Ax\|_\infty\leq C_*\sqrt n\).

The question concerns the sharp worst-case error in simultaneous signed
sums. In matrix form it asks how far a square sign matrix can keep all
sign vectors from the coordinate cube of radius \(C\sqrt n\).

The source also poses \(\limsup_{n\to\infty}d_n>1\) as a related
asymptotic conjecture. It is recorded here as context rather than
counted as another entry.

\subsection{References}\label{references}

\begin{enumerate}
\def\labelenumi{\arabic{enumi}.}
\tightlist
\item
  A. S. Bandeira, A. Kireeva, A. Maillard, and A. Rödder,
  \emph{Randomstrasse101: Open Problems of 2024}, arXiv:2504.20539
  (2025), entry 6, Open Problem 11, Conjecture 12, and the subsequent
  ``Updates'' paragraph.
  \href{https://arxiv.org/html/2504.20539v1}{Paper}.
\item
  A. S. Bandeira, \emph{Did just a couple of deviations suffice all
  along? (problems 10--14)}, Randomstrasse101 (2024), the original
  statement of Open Problem 11.
  \href{https://randomstrasse101.math.ethz.ch/posts/HowManyDeviations/}{Author's
  research post}.
\item
  J. Spencer, \emph{Ten Lectures on the Probabilistic Method}, second
  edition, SIAM (1994), Chapter 10, ``Six Standard Deviations Suffice,''
  pp.~75--80, the classical universal-constant upper bound.
  \href{https://epubs.siam.org/doi/10.1137/1.9781611970074.ch10}{Chapter}.
\end{enumerate}

\subsection{Status check --- 2026-09-10}\label{status-check-2026-09-10}

checked the current arXiv record, still 2504.20539v1 (2025-04-29), the
source's update, and searches ``Spencer sharp constant discrepancy'' and
``Spencer limsup discrepancy 2026''. The update refutes the stronger odd
Sylvester--Hadamard claim in Conjecture 13; that claim is excluded. No
resolution of Open Problem 11 or Conjecture 12 was found. Recent matrix
Spencer results concern a different matrix-valued discrepancy problem.

\textbf{Audit update (2026-09-10):} Rechecked Open Problem 11 and the
source's update, then searched for exact Spencer constants. The refuted
odd Sylvester--Hadamard claim is different. The all-order sharp supremum
warrants extreme rather than hard. This is a bounded literature check,
not a proof that no solution exists.
\end{document}
